\section{2. Related Work}

\subsection{2.1 Sparse Attention}

\textbf{Longformer.} Beltagy et al. (2020) combine a sliding window of radius $W$ 
with a small set of global tokens. Complexity: $\mathcal{O}(L \cdot W)$. 
Content-independent.

\textbf{BigBird.} Zaheer et al. (2020) add random attention and prove Turing 
completeness. Similar to Longformer in practice.

\textbf{Sparse Transformer.} Child et al. (2019) use strided and fixed patterns. 
Pre-dates most modern work.

\textbf{Star-Transformer.} Guo et al. (2019) use a star topology. Rarely used 
today.

\subsection{2.2 Content-Based Sparse Attention}

\textbf{Routing Transformer.} Roy et al. (2021) cluster tokens and restrict 
attention within clusters. Complexity: $\mathcal{O}(L \cdot L/K)$.

\textbf{Reformer.} Kitaev et al. (2020) use LSH to hash similar keys. Probabilistic 
guarantees.

\textbf{MoBA (Mixture of Block Attention).} Lu et al. (2024) introduce block-level 
routing with top-$K$ selection. FUSIONATTN's long path uses MoBA-style 
block selection.

\textbf{NSA (Native Sparse Attention).} DeepSeek (2025) combines local, strided, 
and compressed attention branches.

\subsection{2.3 Hybrid Length-Aware Attention}

\textbf{H-Transformer.} Zhu \& Soricut (2021) adapt pattern to sequence length. 
FUSIONATTN uses a \textbf{discrete dispatch} based on a threshold.

\textbf{Blockwise Parallel Transformer.} Liu \& Abbeel (2023) parallelize block 
attention. Not length-aware.

\subsection{2.4 Grouped Query Attention}

\textbf{GQA.} Ainslie et al. (2023) reduce KV head count. FUSIONATTN preserves 
GQA structure.

\subsection{2.5 Position of Our Work}

\begin{table}[H]
\centering
\begin{tabular}{llll}
\toprule
\textbf{Method} & \textbf{Sparsity} & \textbf{Length-aware} & \textbf{Gate learned} \\
\midrule
Longformer & Fixed & No & No \\
BigBird & Fixed + random & No & No \\
Routing & Content & No & Yes \\
MoBA & Content & No & Yes \\
**FUSIONATTN** & **Content** & **Yes** & **Yes** \\
\bottomrule
\end{tabular}
\end{table}

\textbf{Key distinction.} FUSIONATTN is the only method that **dispatches to 
dense attention for short sequences\textbf{ and }content-aware block sparsity 
for long sequences**.

\subsection{2.6 Comparison Table}

\begin{table}[H]
\centering
\begin{tabular}{lllll}
\toprule
\textbf{Feature} & \textbf{Longformer} & \textbf{Routing} & \textbf{MoBA} & \textbf{FUSIONATTN} \\
\midrule
Short-seq overhead & Yes & Yes & Yes & **None** \\
Long-seq complexity & $O(LW)$ & $O(L^2/K)$ & $O(KL)$ & $O(KL)$ \\
Learned gate & No & Yes & Yes & Yes \\
Multi-arch & Partial & Partial & Partial & Full \\
\bottomrule
\end{tabular}
\end{table}

\subsection{2.7 Not a Compression Method}

FUSIONATTN \textbf{increases} parameter count via LoRA and the gate MLP. It 
should not be confused with compression methods like A3, AASVD, or 
Adaptive SVD.